\documentclass[10pt,a4paper]{article}

\usepackage[T1]{fontenc}
\usepackage[utf8]{inputenc}
\usepackage{lmodern}
\usepackage{amsmath,amssymb,amsthm}
\usepackage{booktabs,array}
\usepackage{geometry}
\usepackage{xcolor}
\usepackage{enumitem}
\usepackage{microtype}
\usepackage{hyperref}
\usepackage{fancyhdr}

\geometry{margin=1in}
\hypersetup{
  colorlinks=true,
  linkcolor=blue!55!black,
  citecolor=blue!55!black,
  urlcolor=blue!55!black,
  pdftitle={On a Roll: A One-Query Universal Forgery against Poulakis--Rolland v2},
  pdfauthor={Martin Feussner}
}
\setlist{nosep,leftmargin=*}
\setlength{\parindent}{1em}
\setlength{\parskip}{0.25em}
\pagestyle{fancy}
\fancyhf{}
\fancyhead[L]{One-query forgery against Poulakis--Rolland v2}
\fancyhead[R]{9 October 2026}
\fancyfoot[C]{\thepage}
\renewcommand{\headrulewidth}{0.3pt}

\newtheorem{theorem}{Theorem}
\newtheorem{observation}{Observation}
\newcommand{\Z}{\mathbb{Z}}
\newcommand{\ord}{\operatorname{ord}}
\newcommand{\Opoint}{\mathcal{O}}
\newcommand{\GH}{G_H}

\title{\bfseries On a Roll:\\
\large A One-Query Universal Forgery against Poulakis--Rolland v2}
\author{Martin Feussner\\
\small Selmer Center, University of Bergen\\
\small \texttt{martin.feussner@uib.no}}
\date{9 October 2026}

\begin{document}
\maketitle

\begin{center}
\fbox{\parbox{0.92\linewidth}{\small\textbf{AI provenance and review status.}
OpenAI Codex generated and selected this paper's title.  It intentionally
chose ``On a Roll'' as wordplay on Rolland.  OpenAI Codex and specialist AI agents performed
the target reconstruction, cryptanalysis, implementations, experiments,
technical and adversarial review, history search, and drafting.  An
independently developed AI implementation reproduced the attack with new
full-parameter key material.  The technical,
adversarial, independent-reproduction, history, and publication-review gates
passed; this report is classified \texttt{VERIFIED\_ATTACK}.  Independent AI
reproduction is not described as human verification.}}
\end{center}
\vspace{0.7em}

\begin{abstract}
We give a one-query universal fresh-message forgery against the second
revision of the composite-order pairing signature of Poulakis and Rolland,
IACR ePrint 2012/134, revision \texttt{20150430:055329}.  Section~4 maps a
message to $H(m)=[\kappa_m]\GH$, where the fixed group generator and the
rejection-sampled scalar $\kappa_m$ are public.  From one ordinary signature
$(s_0,S_0)$, an attacker keeps $s_0$ unchanged and scales $S_0$ by a public
factor involving $\kappa_1/\kappa_0$ and $g^{h(m_1)-h(m_0)}$.  Bilinearity then
turns the source verification equation into the equation for an arbitrary
fresh target.  The attack does not factor $n$, solve a discrete logarithm,
reuse a nonce, exploit malformed encoding, or weaken verifier checks.

We reproduce the attack with the paper's 1023/1024-bit printed factors,
2046-bit composite order, 2048-bit field, supersingular curve, and a complete
modified reduced Tate pairing.  A second implementation generates a new
exact-order subgroup generator, secret exponent, public key, and signing
split at that endpoint.  Both ordinary pairing verifiers accept fresh
forgeries, while changed-message, point, scalar, identity, and non-subgroup
controls reject.  The generic public-scalar transport is known from
Tibouchi's 2011 analysis of insecure BLS hash-to-group constructions; we claim
no novelty for that technique.  Our contribution is its concrete application
and full-parameter validation against Poulakis--Rolland v2.
\end{abstract}

\noindent\fcolorbox{black!45}{black!4}{\parbox{0.94\textwidth}{
\textbf{Result status.} Practical full-parameter design attack: one ordinary
chosen-message query yields fresh-message signatures accepted by the ordinary
verifier.  No exact-target v2 attack was located in the searches completed as
of 9 October 2026.  This is a dated search result, not an absolute priority
claim.}}

\section{Introduction}

Poulakis and Rolland propose a signature intended to retain security unless
both integer factorization and an elliptic-curve discrete logarithm problem
become tractable~\cite{poulakisrolland}.  Its public key combines an RSA-like
composite $n=p_1p_2$, exponentiation in $\Z_n^*$, an order-$n$ elliptic-curve
subgroup, and a bilinear pairing.  The second revision replaced the first
version's signing equations with a randomized split $k+\ell=a$ and a new
pairing verification equation.

The second revision nevertheless specifies a map to the elliptic-curve group
whose output has a publicly known discrete logarithm relative to a fixed
public point.  This representation is enough to transport a valid signature
between messages.  One source query produces a signing response that can be
rescaled for any fresh message for which the advertised affine signature
encoding is defined.

Our contributions are as follows.
\begin{itemize}
  \item We derive a same-response attack that transforms one ordinary v2
  signature into a signature on any distinct encodable target.  Because the
  integer response is copied verbatim, every range check passed by the source
  signature is preserved.
  \item We give a second algebraic form that reconstructs a
  message-independent point and selects a target using a public range bound.
  \item We execute complete 2048-bit-field pairing verifiers in two
  standard-library implementations.  The independent implementation creates
  new full-parameter curve and key material and tests malformed-point and
  non-subgroup controls.
  \item We separate the attack from defects in the paper's printed example,
  from underspecified encodings, and from the already-broken first revision.
  We give a direct repair boundary and explicit prior-art attribution.
\end{itemize}

\paragraph{Attack-history statement.}
The dated search covered the exact title, authors, ePrint identifier and
revision, citations to the target, public code, theses, and terms concerning
forgery, public-scalar maps, hash-to-curve, and BLS-style transport.  It
located the generic technique described below, but no public report applying
it to the exact v2 construction.  Unindexed and private observations cannot be
excluded, so this report makes no absolute first-attack claim.

\paragraph{Technique attribution.}
Boneh, Lynn, and Shacham introduced their pairing-based short-signature
scheme in 2001~\cite{bls}.  Tibouchi's 2011 thesis, Section~3.2,
pages~38--40, analyzes the insecure construction $H(m)=[h(m)]G$ and its
variant $H(m)=[h(m)]Q$~\cite{tibouchi}.  For the latter, one known signature
can be scaled by $h(m)/h(m_0)$ to sign arbitrary messages even when the
discrete logarithm of $Q$ relative to $G$ is hidden.  Poulakis--Rolland v2
adds a scalar response and message-dependent powers of $g$, but those terms
are public and can be included in the same transport.  We claim no novelty
for public-scalar BLS transport.

\section{Exact target and version boundary}

The target is IACR ePrint 2012/134, revision
\texttt{20150430:055329}~\cite{poulakisrolland}.  The audited PDF has SHA-256
digest
\begin{center}
\small\ttfamily
a8b4eae608a28baf86e585fb19056a43881d3c386eae6f7fb6b2cf1d2d49c65f.
\end{center}
The archived author-source gzip has SHA-256 digest
\begin{center}
\small\ttfamily
d67435736b459ef85c705dd726763204ac6400b3c3c9058ac2a2b2e25011ae86.
\end{center}
The revision record states that the first version was broken.  Version~1 is a
materially different system: its public key includes additional points and a
public scalar $r=g^b$, its point response is the static
$S=[g^{ab}]H(m)$, its scalar is $s=bh(m)+a-ab\pmod{\varphi(n)}$, and its
verifier checks two relations.  Version~2 instead uses a fresh split
$k+\ell=a$ and the equations in Section~\ref{sec:scheme}.  Consequently, an
attack or regression test for v1 does not establish a v2 break.  All claims
below are derived from and executed against v2.

\section{The v2 scheme}
\label{sec:scheme}

Key generation chooses $n=p_1p_2$, an elliptic-curve point $P$ of exact order
$n$, a bilinear pairing $e_n$ such that $e_n(P,P)$ has order $n$, a unit
$g\in\Z_n^*$, and a secret $a\in\{1,\ldots,\varphi(n)-1\}$.  It publishes
\[
  Q=[g^a\bmod n]P
\]
and the public key $(P,Q,g,n)$.  The notation $[z]X$ denotes scalar
multiplication in the additive elliptic-curve subgroup.

To sign $m$, the signer chooses positive $k,\ell$ with $k+\ell=a$ and computes
\begin{align}
  s &\equiv k+h(m)+n \pmod{\varphi(n)}, \\
  S &= [g^\ell\bmod n]H(m).
\end{align}
The signature contains $(s,x(S),b(S))$, where the bit selects one of the two
points at the affine $x$-coordinate.  The verifier reconstructs $S$ and accepts
when
\begin{equation}
  e_n([g^s]P,S)
  =e_n([g^{h(m)+n}]Q,H(m)).
  \label{eq:verify}
\end{equation}

\subsection{The public-scalar map}

Section~4 fixes a public generator of the order-$n$ group, called $Q$ in that
section.  To avoid confusing it with the public-key point $Q=[g^a]P$, we call
the map generator $\GH$.  A public KDF is evaluated on $m\mathbin\|T_i$ until
the resulting $|n|$-bit integer is smaller than $n$, and the algorithm outputs
\begin{equation}
  H(m)=[\kappa_m]\GH.
  \label{eq:map}
\end{equation}
Anyone who can evaluate $H$ necessarily computes the accepted
$\kappa_m$.  Hiding the discrete logarithm of $\GH$ relative to $P$ does not
hide ratios between map outputs.

\section{One-query universal forgery}

The adversary first chooses a source message $m_0$ for which
$\gcd(\kappa_0,n)=1$.  This condition is publicly testable before making a
signing query.  Under the stated pseudorandom behavior its failure probability
is approximately $1/p_1+1/p_2$; a nontrivial gcd also reveals a factor of $n$.
The adversary obtains one ordinary signature $\sigma_0=(s_0,S_0)$ on $m_0$.

For any distinct target message $m_1$, define
\begin{align}
  \Delta h &= h(m_1)-h(m_0),\\
  \mu &= \kappa_1\kappa_0^{-1}g^{\Delta h}\pmod n,\\
  s_1 &= s_0,\\
  S_1 &= [\mu]S_0.
  \label{eq:attack}
\end{align}
For negative $\Delta h$, the attacker exponentiates the public inverse of
$g$ modulo $n$.  The output is encoded in the same $(s,x,b)$ format as an
honest signature.

\begin{theorem}
If $\sigma_0$ is an honest signature on $m_0$,
$\gcd(\kappa_0,n)=1$, and $H(m_1)$ is representable in the advertised affine
format, then the signature $(s_1,S_1)$ from Equation~\eqref{eq:attack} is
accepted for $m_1$ by Equation~\eqref{eq:verify}.
\end{theorem}

\begin{proof}
For the signing split $k+\ell=a$, Equations~\eqref{eq:map} and the signing
algorithm give
\[
  S_0=[g^\ell\kappa_0]\GH,
  \qquad
  S_1=[g^{\ell+\Delta h}\kappa_1]\GH.
\]
Because $s_0\equiv k+h(m_0)+n\pmod{\varphi(n)}$ and $g$ is a unit modulo
$n$,
\[
 g^{s_0+\ell+\Delta h}
 \equiv g^{a+h(m_1)+n}\pmod n.
\]
Bilinearity therefore yields
\begin{align*}
 e_n([g^{s_1}]P,S_1)
 &=e_n(P,\GH)^{g^{s_0+\ell+\Delta h}\kappa_1}\\
 &=e_n(P,\GH)^{g^{a+h(m_1)+n}\kappa_1}\\
 &=e_n([g^{h(m_1)+n}]Q,H(m_1)),
\end{align*}
which is the target verification equation.
\end{proof}

The attacker copies $s_0$ exactly, so it automatically preserves every scalar
range restriction satisfied by the oracle response.  It uses no secret key,
factorization, discrete logarithm, $\varphi(n)$, nonce collision, or malformed
point.  After one query, the same computation applies independently to every
fresh encodable target.

\subsection{Range-sampled alternative}

A second form makes the message-independent point explicit:
\[
 c=s_0-h(m_0),\qquad
 U=[\kappa_0^{-1}]S_0.
\]
For a fresh candidate $m_*$, set
\[
 s_*=c+h(m_*),\qquad S_*=[\kappa_*]U.
\]
The hidden modular wrap in $s_0$ does not affect verification because powers
of $g$ are periodic modulo $\varphi(n)$.  If the verifier is expected to apply
the signer-only range $s_*<\varphi(n)$, the attacker samples public targets
until
\[
  1\le s_*\le n-2^{1025}.
\]
The advertised factor bound $p_1,p_2<2^{1024}$ gives
$\varphi(n)=n-p_1-p_2+1>n-2^{1025}$, so this public test guarantees the hidden
range condition.  The same-$s$ attack is stronger because it needs no target
sampling or bound.

\section{Full-parameter implementations}

Both reproductions use the paper's printed factors and therefore the exact
composite-order scale:
\[
 \lvert p_1\rvert=1023,\quad
 \lvert p_2\rvert=1024,\quad
 \lvert n\rvert=2046,\quad
 q=4n-1,\quad \lvert q\rvert=2048.
\]
They implement $E:y^2=x^3+x$ over $\mathbb{F}_q$,
$\mathbb{F}_{q^2}=\mathbb{F}_q[i]/(i^2+1)$, the distortion map
$(x,y)\mapsto(-x,iy)$, a Miller loop at composite order $n$, and final
exponent $(q^2-1)/n=4(q-1)$.  The pairing sanity checks establish
nontriviality, exact order $n$, and bilinearity on independent small scalars.

The paper leaves $h$, the KDF, byte framing, the sign bit, and point encodings
abstract.  Both implementations use domain-separated SHAKE256 and compressed
affine points.  The attack depends only on the public $\kappa_m$ exposed by
the specified map and not on a collision or biased hash output.

\subsection{Primary reconstruction}

The primary implementation uses the published factors, $g=2$, and
$a=2^{256}+2^9+1$.  Its fixed query is
\begin{center}
\texttt{authorized one-query source message}
\end{center}
and its unsigned target is
\begin{center}
\texttt{fresh target message never submitted to the signer}.
\end{center}
The strict verifier decodes the affine point, checks the public scalar range,
checks $[n]S=\Opoint$, and evaluates both sides of
Equation~\eqref{eq:verify}.  An honest signature was accepted, the source
signature was rejected on the target, the same-$s$ forgery was accepted on
the fresh target, and changed-message and altered-scalar controls rejected.

The frozen run took 104.40 seconds wall time and 16,896 KiB peak resident
memory on one worker.  A separate replay of the same frozen source took
105.35 seconds and 17,408 KiB; it repeated all pairing and attack checks.

\subsection{Independent fresh-key reconstruction}

The second implementation independently generates a new exact-order
subgroup generator, secret exponent, public point, and signing split at the
same endpoint.  Its fixed source and preselected target are
\begin{center}
\texttt{single chosen-message signing query}\\
\texttt{fixed fresh target chosen before the attack}.
\end{center}
Neither attack function receives $p_1,p_2,\varphi(n),a,k$, or $\ell$; the
factors remain only in key generation and observer assertions.

The ordinary verifier accepted the honest query, the fixed-target same-$s$
forgery, and the independently exercised range-sampled forgery.  It rejected
the same forgery after changing the message, point, point sign, or scalar, and
rejected both the identity encoding and an order-two non-subgroup point.  The
same-$s$ attack took 2.349 seconds and its pairing verification 18.827
seconds.  The range-sampled form found a target on its first attempt, took
4.540 seconds, and passed verification.  The complete run, including key
generation, three accepted pairing verifications, and all controls, took
156.744 seconds internally (156.83 seconds wall time) and 17,920 KiB peak
resident memory on one worker.

\begin{table}[ht]
\centering
\caption{Recorded full-parameter executions.  Times are wall or internal wall
times on one worker; they are reproducibility measurements rather than
optimized benchmarks.}
\label{tab:runs}
\small
\begin{tabular}{@{}lrrrl@{}}
\toprule
Execution & Full run (s) & attack only (s) & peak KiB & Outcome \\
\midrule
Primary frozen run & 104.40 & included & 16,896 & pass \\
Primary root replay & 105.35 & included & 17,408 & pass \\
Independent fresh key, same-$s$ & 156.744 & 2.349 & 17,920 & pass \\
Independent fresh key, ranged & same run & 4.540 & same run & pass \\
\bottomrule
\end{tabular}
\end{table}

\section{Printed example defects and evidence boundary}

The full-scale endpoint is usable, but the two printed points cannot be taken
verbatim.
\begin{enumerate}
  \item The printed decimal $y(P)$ has an extra digit at zero-based index 587
  and is not on the curve.  Taking the square root determined by the printed
  $x(P)=2^{1500}+2$ repairs the coordinate.
  \item The repaired printed $P$ has order $2n$, with
  $[n]P=(0,0)$ and $[2n]P=\Opoint$, rather than the claimed order $n$.
  \item The printed $Q$ is on the curve and matches $[2^a]P$ for that repaired
  order-$2n$ point, so it has the same subgroup mismatch.  The primary
  implementation doubles both points, obtaining exact-order-$n$ points that
  preserve $Q=[2^a]P$.
  \item Although the text calls both factors 1024-bit primes, the printed
  $p_1$ is 1023 bits and $p_2$ is 1024 bits.
\end{enumerate}
The independent implementation avoids these point defects entirely by
generating a new conforming exact-order point.  It reaches the same attack
result, establishing that the forgery is not caused by repair of the example.

The Section~4 sampler also permits $\kappa_m=0$, while the advertised
signature format cannot encode the resulting identity point by an affine
$x$-coordinate.  Our universal-forgery statement covers every target for
which the scheme's own signature representation is defined.  This negligible
encoding gap neither enables nor assists the attack.

The following hashes bind the executable evidence used for this report:
\begin{center}
\small
\begin{tabular}{@{}p{0.35\linewidth}p{0.58\linewidth}@{}}
\toprule
Artifact & SHA-256 \\
\midrule
Frozen primary source & {\scriptsize\ttfamily dda52b3d3e6e34a0b238ee7034e30f2030f82547e6831d56d6f985251f8fd255} \\
Primary replay output & {\scriptsize\ttfamily eabe8dc9ff3d32cfe1bb5211ee133a5ebd1140df2b76164fa10ed24423581ff2} \\
Independent source before editorial comment fixes & {\scriptsize\ttfamily 6c4c3c675dc474841ea19d09b6c37ac5c663d0c0445f376c548c096fa04718dc} \\
Independent final output & {\scriptsize\ttfamily 9d019f15e0335c350507a5078f3fb5cbcbcb754f73f977b5131b3f8a6ec32d08} \\
Public primary core & {\scriptsize\ttfamily 9b2c4c22ca98a983d7e7df0797f269c32d5cdd838f0956bbb697dc646d4ac921} \\
Public independent core & {\scriptsize\ttfamily 49bf7d87a5904550f2f0cef911c56fa93827aafb83cb7b5c86206825cb5bcb3a} \\
\bottomrule
\end{tabular}
\end{center}
The public copies change only explanatory comments and docstrings: they
correct the factor-size wording, identify $\GH$ rather than $P$ as the
Section~4 map base used by the implementation, and remove development
provenance.  Executable logic is unchanged.

\section{Repair and limitations}

The direct cause is Equation~\eqref{eq:map}: public map evaluation exposes a
discrete logarithm of every output with respect to the same public base.
Changing that base from $P$ to an unrelated public point does not help, as
Tibouchi's BLS analysis already shows; the attack needs scalar ratios, not the
base's discrete logarithm.

A repair should replace Section~4 with a standard hash-to-group construction
whose outputs do not carry publicly known discrete logarithms relative to a
fixed base.  The signing and verification equations then require a new
security proof under a precisely stated hash-to-group model.  Implementations
should also specify byte encodings, reject the identity, validate subgroup
membership, and define the public scalar range.  These validation rules are
sound engineering requirements, but the reproduced strict verifiers show
that they do not stop the present transport.

This report demonstrates a design attack at one concrete endpoint because v2
publishes one full numerical example.  Its proof is parameter-independent for
every well-formed key satisfying the stated equations.  The work does not
claim secret-key recovery, factorization, or a break of the underlying
pairing, discrete-logarithm, or factoring assumptions.  The stronger direct
forgery already violates existential unforgeability under chosen-message
attack.

\section{Reproducibility}

The public package uses only Python's standard library.  From the attack
directory, run
\begin{verbatim}
python3 reproduce.py
python3 reproduce-all-rows.py
\end{verbatim}
or execute \path{reproducer/reproduce.sh}.  The first command runs the primary
published-example reconstruction; the second independently generates new key
material at the sole published endpoint.  Canonical outputs are stored in
\path{reference-output.json} and
\path{all-rows-reference-output.json}.  Timing fields vary by machine.

The implementation is intentionally direct rather than optimized.  It uses
affine curve arithmetic and pure-Python extension-field operations so every
group and pairing step remains inspectable.  Peak memory stayed below 18 MiB
in the recorded runs; each command occupies one CPU.

\section{Conclusion}

Poulakis--Rolland v2 publishes the scalar representation of every
hash-to-group output.  One valid signature therefore supplies a point that
can be rescaled, together with public message-dependent powers of $g$, into a
valid signature for any fresh encodable target.  Keeping the integer response
unchanged makes the attack compatible with strict range, decoding, and
subgroup checks.  Two full-parameter modified Tate-pairing implementations,
including one with new key material, accept the forgery and reject the
negative controls.  Repair requires a new hash-to-group construction and a
new proof for the resulting scheme.

\section*{Acknowledgements}

The computations were performed on the Norwegian Research and Education Cloud (NREC), using resources provided by the University of Bergen and the University of Oslo.

\begin{thebibliography}{10}

\bibitem{poulakisrolland}
D.~Poulakis and R.~Rolland.
\newblock A digital signature scheme based on two hard problems.
\newblock Cryptology ePrint Archive, Paper 2012/134, revision
  \texttt{20150430:055329}, 2015.
\newblock \url{https://eprint.iacr.org/2012/134}.

\bibitem{tibouchi}
M.~Tibouchi.
\newblock \emph{Hachage vers les courbes elliptiques et cryptanalyse de
  sch\'emas RSA} (Hashing to elliptic curves and cryptanalysis of RSA-based
  schemes).
\newblock Ph.D. thesis, Universit\'e Paris Diderot and University of
  Luxembourg, 2011.  Section~3.2, pp.~38--40.
\newblock \url{https://hdl.handle.net/10993/15584}.

\bibitem{bls}
D.~Boneh, B.~Lynn, and H.~Shacham.
\newblock Short signatures from the Weil pairing.
\newblock \emph{Journal of Cryptology}, 17(4):297--319, 2004.  Preliminary
  version in ASIACRYPT 2001.
\newblock \url{https://doi.org/10.1007/s00145-004-0314-9}.

\end{thebibliography}

\end{document}
